\section{Bevezetés}

A Parkinson-kór a központi idegrendszer progresszív megbetegedése, amely elsősorban a
mozgásszabályozást érinti. A betegség gyakori motoros tünetei közé tartozik a remegés,
az izommerevség, a mozgáslassulás (bradykinesia), valamint az egyensúly és a koordináció romlása. \cite{author2025physiotherapy}
Ezek a tünetek a betegség előrehaladtával egyre nagyobb mértékben befolyásolhatják a
beteg járását, testtartását, mindennapi aktivitását és önállóságát.

A gyógyszeres kezelés mellett fontos szerepet kap a mozgásterápia és a fizioterápia,
amelyek célja a mozgásfunkciók fenntartása, a járás és az egyensúly javítása, valamint az
életminőség támogatása.\cite{osborne2022physical} rehabilitáció általában hosszabb távú folyamat,
amelynek eredményességét jelentősen befolyásolhatja a gyakorlatok rendszeressége
és helyes végrehajtása.


A terápiás alkalmakon kívül végzett gyakorlatok esetében kihívást jelenthet a végrehajtás
követése és értékelése. A szakember számára gyakran korlátozott mennyiségű információ áll
rendelkezésre arról, hogy a beteg a gyakorlatokat milyen módon és milyen minőségben végezte el.
Emiatt a rehabilitációban egyre fontosabbá válnak azok a digitális megoldások,
amelyek a gyakorlatokhoz kapcsolódó adatokat mérhető és később is elemezhető formában rögzítik.
\cite{rossi2022parkinson}


A mobilalkalmazások és távrehabilitációs rendszerek lehetőséget adnak arra, hogy a beteg a
terápiás alkalmakon kívül is irányított formában végezzen feladatokat.
Egy ilyen funkciókat ellátó alkalmazás nemcsak a feladatok megjelenítésében segíthet,
hanem a mérési folyamat egységesítésében és az adatok rendszerezésében is.
Ez különösen olyan esetekben hasznos, amikor a hosszabb távú követés vagy a
gyakorlatok összehasonlíthatósága fontos szempont. \cite{chen2024application}


A digitális rehabilitáció egyik lényeges eleme az objektív mérési adatok gyűjtése.
A viselhető szenzorok előnye, hogy a mozgás közben, a felhasználót viszonylag kevéssé
zavarva képesek adatokat szolgáltatni. \cite{paci2020design} Parkinson-kór esetén különösen hasznosak lehetnek azok
a szenzorok, amelyek a járáshoz, egyensúlyhoz vagy testsúlyáthelyezéshez kapcsolódó információkat
mérik. A lábon viselt eszközök ebből a szempontból kiemelt jelentőségűek, mivel közvetlenül a
talajkontaktusról, a terhelési mintázatról és a láb mozgásáról adhatnak információt.


A Sensoria okoszokni ilyen típusú viselhető mérőeszköz, amely a talp különböző részein
jelentkező nyomásviszonyokról és a láb mozgásáról szolgáltat adatokat. \cite{sensoria2016walk}
Ez a mérési megoldás jól illeszkedik egy mobilalkalmazással támogatott rehabilitációs rendszerhez,
mivel nem igényel nagyméretű, helyhez kötött laboratóriumi eszközöket.
Az így gyűjtött adatok alapot adhatnak a gyakorlatok későbbi áttekintéséhez, összehasonlításához
és további elemzéséhez.


A jelen dolgozat egy ilyen, Sensoria okoszoknival együttműködő mobilalkalmazás fejlesztését
mutatja be. A következő fejezetekben az alkalmazás felépítése, fő komponensei, az adatrögzítés
folyamata, valamint az eredmények megjelenítésének módja kerül ismertetésre.
















